\section{Methodology}

\subsection{Timing}

All kernel timing uses CUDA events, not host wall clock, so the measurement sees
device time and is not polluted by launch overhead or host scheduling. To keep
the per launch overhead from dominating on the cheap kernels, I place many
launches between a single pair of events and divide by the launch count rather
than synchronizing once per launch. Before any timed run I discard a batch of
warmup iterations, which lets the boost clock ramp and the caches warm so the
timed region measures steady state rather than the cold transient.

I then time at least ten batches and report the mean, median, and standard
deviation. This card boosts its clock opportunistically and shares silicon with
the display, so run to run variance is real and worth showing rather than hiding
behind a single number. The standard deviation column in every timing table is
there for exactly that reason.

\subsection{Counting floating point operations}

Each kernel's FLOP count comes from its actual arithmetic, not an estimate. A
GEMM of dimensions $M \times N \times K$ does $2MNK$ floating point operations,
counting each multiply and add separately. A GEMV of an $M \times N$ matrix does
$2MN$. A SAXPY over $N$ elements does $2N$, one multiply and one add per element.
A transpose does zero floating point work, which is the whole point of including
it: it isolates the memory system with no arithmetic to hide behind.

\subsection{Counting bytes}

I report two byte counts wherever both are available. The theoretical count is
the minimum traffic a perfect implementation would generate: read each input
once, write each output once, in the kernel's precision. The measured count is
the DRAM read and write bytes Nsight Compute observed. Arithmetic intensity is
computed from the measured bytes where the profiler ran, and falls back to the
theoretical bytes elsewhere, always labeled so the two are never confused.

\subsection{Theoretical versus empirical peak}

The theoretical ceilings come from the device properties queried at runtime:
the FP32 peak from the SM count, the FP32 lanes per SM, and the boost clock,
counting a fused multiply add as two operations; the bandwidth peak from the bus
width and the effective memory transfer rate. The empirical ceilings come from
pushing the hardware directly: a register resident FMA loop for compute, and a
large device to device copy for bandwidth. A card almost never reaches its
theoretical peak, and reporting the measured ceiling next to the theoretical one
is what keeps the roofline honest about what this specific machine delivers.
